\section{证明程序}
\begin{enumerate}
  \item 证据的保全。
        \begin{enumerate}
          \item 诉讼中。证据可能灭失 \hz 日后难以取得的。
                \begin{enumerate}
                  \item 当事人申请 \hz 法院依职权。
                  \item 举证期届满前，受理案件的法院管辖。
                  \item 可能给他人造成损失的应当责令提供担保（\zd{担保你瞎申请保全}）。
                \end{enumerate}
          \item 起诉前 \hz 仲裁前。
                \begin{enumerate}
                  \item 证据所在地、被申请人所在地、将来对案件有管辖权的法院（结合案件）。
                  \item 48 小时内下裁定，30 天内起诉申请，否则法院应当裁定解除保全。
                  \item 应当提供担保（只听你一面之词，风险较大）。
                \end{enumerate}
        \end{enumerate}
  \item 举证。
        \begin{enumerate}
          \item 协商 \hz 法院确定的\zd{举证期限}不得少于 15 天。
          \item 延长。举证期限届满前，书面申请。
          \item \zd{逾期举证}的后果。主观恶性大小 + 证据重要程度。
                \begin{enumerate}
                  \item 主观恶性特别大的、故意 \hz 重大过失的 \yjt 不采纳。
                  \item 故意、重大过失但\zd{证据很重要}（与基本事实相关）\yjt 采纳 + 训诫 + 罚款。
                  \item 不是故意不是重大过失\yjt 采纳 + 训诚。
                  \item 客观原因，对方无异议 \yjt 视为未逾期，正常采纳，无需训诫、罚款。
                \end{enumerate}
        \end{enumerate}
  \item 调查收集证据。
        \begin{enumerate}
          \item 法院依职权。\zd{三大类，五种情形}：
                \begin{enumerate}
                  \item 国家、社会、第三人利益的 \yjt 不会破坏双方平等。
                  \item \zd{身份关系} \yjt 涉及公序良俗。
                  \item 追加当事人、中止、终结诉讼、回避等\zd{程序性事项} \yjt 不涉及双方实体权利。
                  \item \zd{不需要质证}。跟当事人关系不大，没有质证积极性；法官不能成为质证主体。
                \end{enumerate}

                \begin{kk}{法院查明的事实与当事人主张的事实不一致}
                  \begin{enumerate}
                    \item 禁止突袭裁判。
                    \item \zd{归纳为争议焦点}，当事人可以变更诉讼请求，法院可以重新指定举证期限。
                  \end{enumerate}
                \end{kk}
          \item 原则依申请。举证期限内，书面申请，正常质证。
        \end{enumerate}
  \item 证据必须经过质证，必须在\zd{公开开庭}时进行（涉密、隐私除外）。
  \item \zd{证据裁判原则/主义}。应当以证据依法裁判，认定事实靠证据。
        \begin{enumerate}
          \item \yi \zd{证据能力/资格} \yjt 法律问题 \yjt 法律规定 \er \zd{证明力} \yjt 常识/经验问题 \yjt 法律不规定\zd{自由心证}。
                \footnote{例外：只有被告人供述，没有其他证据证明的，不准定罪量刑。}。
          \item 纠问制（身兼控诉与审判）\yjt 法定证据制度\footnote{自由心证对立面，对每一个证据的证明力，予以明文规定。典型如\zd{口供}，赋予了极高的证明力。} \yjt 自由心证。
          \item 心证公开。推导过程写到判决书中去。约束自由心证。
          \item 非法方法：窃取、窃听、窃录。\zd{遏制非法取证行为}。
          \item \zd{偷拍偷录}不叫非法方法（法律没有禁止）。具体问题具体分析，内容是否侵犯他人利益，违背公序良俗。
        \end{enumerate}
  \item \zd{人世间哪来那么多当然的事情}，如果要蒙肯定蒙否。
\end{enumerate}
